% !PRINTED 111
\weberpage{111}
%<<CONT

\begin{equation}
 y^2+z^2=\sqrt[n]{b^2+c^2},
\tag{8}
\end{equation}
wodurch der absolute Werth von $x$ bestimmt ist. Es ergiebt sich ferner aus (6) und (7)
\begin{equation}
 \varphi(y,z)={ (y-iz)^n(b+ic)-(y+iz)^n(b-ic) \over 2i}=0.
\tag{9}
\end{equation}
Nun ist $\varphi(y,z)$ eine ganze homogene Function $n$ten Grades der beiden Veränderlichen $y,z$, und zwar mit reellen Coëfficienten (da die Vertauschung von $i$ mit $-i$ in $\varphi(y,z)$ nichts ändert), und wenn wir nach absteigenden Potenzen von $y$ ordnen, so ist das erste Glied $cy^n$, also nach Voraussetzung von Null verschieden. Setzen wir nun
\[
 y=\lambda z,
\]
so erhalten wir aus (9)
\begin{equation}
 \varphi(\lambda,1)=0,
\tag{10}
\end{equation}
also eine Gleichung $n$ten Grades für $\lambda$, die nach X. gewiss eine Wurzel hat. Ist $\lambda$ bestimmt, so folgt aus (8)
\begin{equation}
 z=\sqrt{\frac{\sqrt[n]{b^2+c^2}}{1+\lambda^2}},\qquad
 y=\lambda\sqrt{\frac{\sqrt[n]{b^2+c^2}}{1+\lambda^2}}.
\tag{11}
\end{equation}
Hierdurch sind die Gleichungen (8), (9) thatsächlich befriedigt; aus (9) aber folgt
\[
 \frac{(y+iz)^n}{b+ic}=\frac{(y-iz)^n}{b-ic},
\]
und aus (8), dass der gemeinsame Werth dieser beiden Brüche $\pm 1$ ist. Wir haben daher
\[
 (y+iz)^n=\pm(b+ci),\qquad (y-iz)^n=\pm(b-ci),
\]
und es kann also das Vorzeichen der Quadratwurzel in (11) so bestimmt werden, dass die Gleichungen (6) und (7) erfüllt sind.

\begin{center}
{\large \S.\ 33.}

\vspace{0.5em}
{\large\gesperrt{Lösung reiner Gleichungen durch trigonometrische Functionen.}}
\end{center}

Weit vollständiger und einfacher kann die Auflösbarkeit einer reinen Gleichung dargethan werden, wenn man die trigonometrischen Functionen und ihre einfachsten Eigenschaften als bekannt voraussetzt. Freilich sind diese Functionen der
%CONT>>
